(1) 光程差 \(\Delta \) 引起时间延迟 \(\tau=\Delta/c\)，对于单频光的两束等光强 \(I_{0}\)，输出总强度可写为 \(I_\omega(\Delta)=2I_0[1+\cos(\omega\tau)]\)。实际高斯谱的不同频率分量在探测器的慢响应下互不相干，须对频率功率而不是不同频率振幅直接求和。

(2) 取归一化谱
\[
s(\omega)=\frac1{\sqrt{2\pi}\sigma}
e^{-(\omega-\omega_0)^2/(2\sigma^2)}
\]
并在 \(\sigma\ll\omega_0\) 时将下限近似延伸至负无穷。高斯特征函数满足
\[
\int_{-\infty}^\infty s(\omega)e^{i\omega\tau}d\omega
=e^{i\omega_0\tau-\sigma^2\tau^2/2}.
\]
因此
\[
I(\Delta)=2I_0\left[1+
e^{-\sigma^2\Delta^2/(2c^2)}
\cos(\omega_0\Delta/c)\right].
\]
局部干涉可见度包络为
\[
\mathcal V(\Delta)=e^{-\sigma^2\Delta^2/(2c^2)}.
\]

(3) 按 \(\mathcal V(\ell_{\rm coh})=e^{-1}\)，求得
\[
\ell_{\rm coh}=\frac{\sqrt2\,c}{\sigma}.
\]
对应相干时间为 \(\tau_{\rm coh}=\sqrt2/\sigma\)。相干长度大表示光源频谱窄；该数值系数依赖采用角频率标准差及 1/e 包络的约定。

(4) 实验上缓慢扫描 \(\Delta \) 可获得快速载波振荡的周期 \(2\pi c/\omega_0\) 和慢变高斯包络的宽度，通过 \(\ln\mathcal V=-\sigma^2\Delta^2/(2c^2)\) 对 \(\Delta ^{2}\) 作直线拟合即可估计 \(\sigma \)。若 \(\sigma \to 0^{+}\)，包络趋1，任意有限光程差仍有完全对比度，退化为单色光；有限\(\sigma \)导致可见度下降，并不意味着两臂本身存在吸收。
